\begin{UPSTIinfor}{Notation utilisée dans ce TD : le pseudo-code}
    \label{info:pseudocode}
    Avant d'écrire du code C, il est utile de décrire \textbf{l'idée} d'un programme avec une notation plus proche du français, indépendante de tout langage. C'est ce qu'on appelle du \textbf{pseudo-code}. Il joue le même rôle que les blocs de Scratch : décrire un enchaînement d'actions et des choix, sans se soucier de la syntaxe exacte d'un langage donné.

    Dans ce TD, on utilisera les conventions suivantes.

    \begin{description}
        \item[Affectation :] \verb|variable <- expression| se lit \og la variable prend la valeur de l'expression \fg{}. C'est l'équivalent du bloc Scratch \og mettre \texttt{variable} à ... \fg{}.
        \item[Structure conditionnelle :]
            \begin{verbatim}
SI condition ALORS
    instruction(s)
SINON SI condition ALORS
    instruction(s)
SINON
    instruction(s)
FIN SI
            \end{verbatim}
            C'est l'équivalent de l'empilement de blocs Scratch \og si ... alors ... sinon \fg.
              Les blocs \texttt{SINON SI} et \texttt{SINON} sont facultatifs, tout comme en C.
        \item[Comparateurs :] on garde les mêmes symboles qu'en C : \verb|==|, \verb|!=|, \verb|<|, \verb|>|, \verb|<=|, \verb|>=|.
        \item[Opérateurs logiques :] on les écrit en toutes lettres, pour ne pas les confondre avec ceux du C : \textbf{ET}, \textbf{OU}, \textbf{NON}. (\texttt{ET} correspond à \texttt{\&\&}, \texttt{OU} à \texttt{||}.)
        \item[Affichage :] \verb|AFFICHER "un texte"| ou \verb|AFFICHER variable| (équivalent d'un \texttt{printf}).
    \end{description}

    \textbf{À retenir :} le pseudo-code décrit \emph{ce que fait} le programme ; le C décrit \emph{comment} l'écrire pour que l'ordinateur l'exécute. Un même pseudo-code peut se traduire en C de plusieurs façons correctes.
\end{UPSTIinfor}

\begin{UPSTIexemples}[Une même idée en trois écritures]
    \label{ex:trois-ecritures}
    Voici comment la même idée --- \og retenir un score, puis dire si la partie est gagnée \fg{} --- s'exprime dans les trois notations que l'on utilisera.

    \bigskip
    \noindent
    \begin{tabular}{|l|l|l|}
        \hline
        \textbf{Scratch}                                & \textbf{Pseudo-code}       & \textbf{C}             \\
        \hline
        bloc \og mettre \texttt{score} à \texttt{0} \fg{} & \verb|score <- 0|         & \verb|int score = 0;| \\
        \hline
    \end{tabular}

    \bigskip
    \noindent
    Et pour le choix \og si le score dépasse 10, la partie est gagnée, sinon elle est perdue \fg{}, obtenu en empilant des blocs \og si ... alors ... sinon \fg{} en Scratch :

    \medskip
    \noindent
    \begin{minipage}[t]{0.32\linewidth}
        \textbf{Pseudo-code}
        \begin{verbatim}
SI score > 10
ALORS
  AFFICHER "Gagne"
SINON
  AFFICHER "Perdu"
FIN SI
        \end{verbatim}
    \end{minipage}\hfill
    \begin{minipage}[t]{0.32\linewidth}
        \textbf{C}
        \begin{lstlisting}[language=C]
if (score > 10)
{
  printf("Gagne\n");
}
else
{
  printf("Perdu\n");
}
        \end{lstlisting}
    \end{minipage}\hfill
    \begin{minipage}[t]{0.32\linewidth}
        \textbf{Scratch (description)}

        Bloc \og si \texttt{score} $>$ \texttt{10} alors \fg{} contenant un bloc \og dire "Gagne" \fg, avec une branche \og sinon \fg{} contenant un bloc \og dire "Perdu" \fg.
    \end{minipage}
\end{UPSTIexemples}
